\subsection{Eficiencia}
\label{sec:eficiencia}

Entre dos algoritmos correctos que resuelven un mismo problema, naturalmente se prefiere el que menos recursos utiliza: el más \concept{eficiente}.

Usualmente, el recurso más preciado es el tiempo, por lo cual el análisis de eficiencia suele estar enfocado en el tiempo de ejecución; sin embargo, se pueden seguir principios similares para analizar otros recursos, como el uso de memoria.

\subsection{Complejidad temporal}

El tiempo de ejecución de una implementación de un algoritmo depende de:
\begin{enumerate}
  \item La máquina en la que se ejecuta.
  \item El lenguaje de programación en el que se implementa.
  \item El número de operaciones que realiza el algoritmo.
  \item Los datos que recibe el programa como entrada.
\end{enumerate}

Al analizar un algoritmo, se pretende sacar conclusiones generales; por ende, no basta con implementar el algoritmo en una máquina y lenguaje específicos para medir su tiempo de ejecución. Por el contrario, se intenta estimar el número de operaciones elementales que realiza el algoritmo. Para eso, en lugar de ejecutar el algoritmo en una máquina real, se analiza el algoritmo en el marco de un \concept{modelo de computación} que abstrae la máquina física y define cuáles operaciones son elementales.

Usualmente, la cantidad de operaciones depende de qué tantos datos el algoritmo recibe como entrada. Por ello, la cantidad se expresa como una función matemática del tamaño de la entrada. Esa función se denomina \concept{costo computacional temporal}.

Se puede analizar el crecimiento asintótico de esa función para saber cómo escala el costo del algoritmo a medida que aumenta el tamaño de la entrada. El resultado de ese análisis es la \concept{complejidad temporal} del algoritmo y suele ser el objetivo último del análisis de eficiencia.

De forma análoga, el modelo de computación define cuánto espacio en memoria ocupan las operaciones, de forma que se puede definir una \concept{función de costo computacional espacial} y analizarla asintóticamente en función de la entrada para conocer la \concept{complejidad espacial}.


